\documentclass[a4paper,
fontsize=11pt,
%headings=small,
oneside,
numbers=noperiodatend,
parskip=half-,
bibliography=totoc,
final
]{scrartcl}

\usepackage[babel]{csquotes}
\usepackage{synttree}
\usepackage{graphicx}
\setkeys{Gin}{width=.4\textwidth} %default pics size

\graphicspath{{./plots/}}
\usepackage[ngerman]{babel}
\usepackage[T1]{fontenc}
%\usepackage{amsmath}
\usepackage[utf8x]{inputenc}
\usepackage [hyphens]{url}
\usepackage{booktabs} 
\usepackage[left=2.4cm,right=2.4cm,top=2.3cm,bottom=2cm,includeheadfoot]{geometry}
\usepackage[labelformat=empty]{caption} % option 'labelformat=empty]' to surpress adding "Abbildung 1:" or "Figure 1" before each caption / use parameter '\captionsetup{labelformat=empty}' instead to change this for just one caption
\usepackage{eurosym}
\usepackage{multirow}
\usepackage[ngerman]{varioref}
\setcapindent{1em}
\renewcommand{\labelitemi}{--}
\usepackage{paralist}
\usepackage{pdfpages}
\usepackage{lscape}
\usepackage{float}
\usepackage{acronym}
\usepackage{eurosym}
\usepackage{longtable,lscape}
\usepackage{mathpazo}
\usepackage[normalem]{ulem} %emphasize weiterhin kursiv
\usepackage[flushmargin,ragged]{footmisc} % left align footnote
\usepackage{ccicons} 
\setcapindent{0pt} % no indentation in captions
\usepackage{xurl} % Breaks URLs

%%%% fancy LIBREAS URL color 
\usepackage{xcolor}
\definecolor{libreas}{RGB}{112,0,0}

\usepackage{listings}

\urlstyle{same}  % don't use monospace font for urls

\usepackage[fleqn]{amsmath}

%adjust fontsize for part

\usepackage{sectsty}
\partfont{\large}

%Das BibTeX-Zeichen mit \BibTeX setzen:
\def\symbol#1{\char #1\relax}
\def\bsl{{\tt\symbol{'134}}}
\def\BibTeX{{\rm B\kern-.05em{\sc i\kern-.025em b}\kern-.08em
    T\kern-.1667em\lower.7ex\hbox{E}\kern-.125emX}}

\usepackage{fancyhdr}
\fancyhf{}
\pagestyle{fancyplain}
\fancyhead[R]{\thepage}

% make sure bookmarks are created eventough sections are not numbered!
% uncommend if sections are numbered (bookmarks created by default)
\makeatletter
\renewcommand\@seccntformat[1]{}
\makeatother

% typo setup
\clubpenalty = 10000
\widowpenalty = 10000
\displaywidowpenalty = 10000

\usepackage{hyperxmp}
\usepackage[colorlinks, linkcolor=black,citecolor=black, urlcolor=libreas,
breaklinks= true,bookmarks=true,bookmarksopen=true]{hyperref}
\usepackage{breakurl}

%meta
%meta

\fancyhead[L]{A. Roth\\ %author
LIBREAS. Library Ideas, 48 (2026). % journal, issue, volume.
\href{https://doi.org/10.18452/37904}{\color{black}https://doi.org/10.18452/37904}
{}} % doi 
\fancyhead[R]{\thepage} %page number
\fancyfoot[L] {\ccLogo \ccAttribution\ \href{https://creativecommons.org/licenses/by/4.0/}{\color{black}Creative Commons BY 4.0}}  %licence
\fancyfoot[R] {ISSN: 1860-7950}

\title{\LARGE{Qualitative Interviews und Forschungsethik}}% title
\author{Alexandra Roth} % author

\setcounter{page}{1}

\hypersetup{%
      pdftitle={Qualitative Interviews und Forschungsethik},
      pdfauthor={Alexandra Roth},
      pdfcopyright={CC BY 4.0 International},
      pdfsubject={LIBREAS. Library Ideas, 48 (2026)},
      pdfkeywords={Qualitative Interviews, Methode, Ethik, qualitative interviews, methodology, ethics},
      pdflicenseurl={https://creativecommons.org/licenses/by/4.0/},
      pdfcontacturl={http://libreas.eu},
      pdfurl={https://doi.org/10.18452/37904},
      pdfdoi={10.18452/37904},
      pdflang={de},
      pdfmetalang={de}
     }



\date{}
\begin{document}

\maketitle
\thispagestyle{fancyplain} 

%abstracts
\begin{abstract}
\noindent
\textbf{Kurzfassung}: Mit qualitativen Interviews lassen sich
unterschiedliche Fragestellungen bearbeiten. Die Methode konfrontiert
Forschende allerdings mit einer Reihe ethischer Herausforderungen. Der
Beitrag beschreibt Wertkonflikte, die sich bei Untersuchungen, die
mithilfe qualitativer Interviews durchgeführt werden, ergeben können,
und zeigt mögliche Lösungswege auf.

\begin{center}\rule{0.5\linewidth}{0.5pt}\end{center}

\noindent\textbf{Abstract}: Qualitative interviews can be used to address a wide
range of research questions. However, the method confronts researchers
with a number of ethical challenges. The article describes potential
conflicts of values that can arise in studies using qualitative
interviews and suggests possible solutions.
\end{abstract}

\vspace{0.5 cm}

%body
Der Einsatz qualitativer Interviews bietet sich vor allem dann an, wenn
das Erkenntnisinteresse besteht, \enquote{individuelle Sichtweisen zu
erschließen und sozialen Sinn zu rekonstruieren. Besonders fruchtbar ist
die qualitative Befragung daher für Untersuchungsfragen, zu denen nur
das Individuum Auskunft geben kann: zur Biografie, zu Werten,
Einstellungen, Images und Gewohnheiten} (Werner 2013, S. 129 f.). Der
Einsatz von qualitativen Interviews ist damit auch im Bibliothekskontext
sinnvoll, so etwa, wenn man Genaueres über die Motive der
Bibliotheksnutzer*innen beziehungsweise der Nicht-Nutzer*innen erfahren
möchte, oder wenn man sich, beispielsweise im Rahmen einer studentischen
Arbeit, einem Thema explorativ-hypothesengene\-rierend nähern will.

Wie bei allen anderen wissenschaftlichen Methoden ist bei qualitativen
Befragungen darauf zu achten, dass den Proband*innen aus einer
Untersuchungsteilnahme keine Nachteile erwachsen (vgl.
Przyborski/Wohlrab-Sahr 2021, S. 45). Was in der Theorie recht simpel
klingt, kann sich in der Praxis mitunter als schwierig erweisen.
Problematisch wird es vor allem dann, wenn moralische Prinzipien mit
beruflichen Werten oder mit wissenschaftlichen Gütekriterien
konfligieren. Im Folgenden werden beispielhaft Dilemmata beschrieben,
die qualitative Interviews mit sich bringen können.

\section{Konflikt: Validität und
Transparenz}\label{konflikt-validituxe4t-und-transparenz}

Ein forschungsethisches Dilemma kann sich zunächst durch den Konflikt
zwischen der Forderung, gegenüber den Proband*innen transparent zu
kommunizieren, einerseits und dem Anspruch, valide Ergebnisse zu
erzielen, andererseits ergeben. Transparenz bedeutet in diesem
Zusammenhang, dass Forschende potenzielle Interviewpartner*innen im
Rahmen einer informierten Einwilligung noch vor der eigentlichen
Erhebung über ihre Rechte, die Ziele und die Inhalte der Untersuchung
sowie über etwaige negative Konsequenzen, die sich aus einer Teilnahme
an der Befragung ergeben, informieren müssen (vgl. Misoch 2015, S.
15--23; Seidman 2006, S. 61--69). Auf dieser Basis sollen die
angefragten Personen (beziehungsweise bei Kindern deren
Erziehungsberechtigte) frei und selbstbestimmt entscheiden können, ob
eine Interviewteilnahme erfolgt (vgl. Friedrichs 2022, S. 349; Seidman
2006, S. 61). Eine umfassende Information der Teilnehmenden kann jedoch
dazu führen, dass diese bei der eigentlichen Erhebung sozial erwünschte
Antworten geben (vgl. Misoch 2015, S. 15--23). Dieses Dilemma lässt sich
nur schwer lösen, kann jedoch \enquote{gemildert werden, indem das Ziel
und die Methode der Studie zwar transparent gemacht werden, ohne dabei
aber zu sehr ins Detail zu gehen, sodass die zu Befragenden hinsichtlich
ihres Interviewverhaltens nicht beeinflusst werden} (Misoch 2015, S.
19).

\section{Konflikt: Probandenschutz versus intersubjektive
Nachvollziehbarkeit und
Validität}\label{konflikt-probandenschutz-versus-intersubjektive-nachvollziehbarkeit-und-validituxe4t}

Ein weiteres forschungsethisches Dilemma kann sich ergeben, wenn der
Schutz der Proband*in\-nen mit den wissenschaftlichen Gütekriterien
intersubjektive Nachvollziehbarkeit und Validität konfligiert. So sind
Zusatzinformationen zu Proband*innen, etwa zu ihrem Alter oder ihrem
Arbeitsort, gerade im qualitativen Paradigma hilfreich, um das Gesagte
besser zu kontextualisieren und damit nachvollziehbarer zu machen (vgl.
Roth 2024, S. 190). Auf der anderen Seite könnten derlei Daten die
Re-Identifizierung von Befragungsteilnehmer*innen ermöglichen. Dieses
Risiko besteht vor allem dann, wenn die Population, aus der das Sample
gewählt wird, von vornherein nur einen kleinen Personenkreis umfasst,
weil dann nicht davon auszugehen ist, dass abgefragte Merkmale auf mehr
als eine Person zutreffen. Die Gefahr der Re-Identifizierung kann also
im Bibliothekskontext zum Beispiel bei Mitarbeiterbefragungen in
kleineren Einrichtungen bestehen, aber auch in größeren Systemen, wenn
sehr viele Zusatzinformationen zu den Proband*innen erhoben werden. Hier
gilt es genau abzuwägen, wie viele Daten tatsächlich dem Erkenntnisziel
dienen und im Zweifelsfall großzügig zu anonymisieren.

\section{Konflikt: Fairness versus
Effizienz}\label{konflikt-fairness-versus-effizienz}

Der Schutz der Proband*innen vor Schaden ist letztlich auch ein Zeichen
von Respekt. Ein respektvoller Umgang gegenüber den
Interviewpartner*innen beinhaltet zudem, dass manipulative oder
despektierliche Fragen vermieden werden sollten. Vielmehr ist den
Proband*innen sachlich zu begegnen. Außerdem sollte man es unterlassen,
ihnen Antworten in den Mund zu legen oder sie vor allem auf bestimmte
Rollen, etwa die der Frau, der*s Schüler*in et cetera zu reduzieren. Das
sorgfältige Prüfen und Evaluieren der Frageformulierungen, um solche
Fauxpas zu vermeiden, ist daher notwendig, geht aber zu Lasten der
Effizienz. Gleiches gilt für eine sorgsame Transkription. Diese ist
deshalb so wichtig, weil Abweichungen zu Fehlinterpretationen führen
können (vgl. Misoch 2015, S. 15--23). Wie viel Zeit in die Vor- und
Nachbereitung von Interviews investiert werden sollte, lässt sich zwar
nicht pauschal festlegen. Ein Pretest im Vorfeld der Erhebung und eine
Prüfung der Transkripte auf Korrektheit sind auf jeden Fall angebracht.

\section{Konflikt: Probandenschutz versus
Transparenz}\label{konflikt-probandenschutz-versus-transparenz}

Das Führen von Interviews erfolgt in der Regel in der Absicht, die
Forschungsdaten im Nachhinein einem größeren Kreis zugänglich zu machen.
Damit den Proband*innen daraus keine Nachteile erwachsen, müssen
Forschungsdaten so anonymisiert werden, dass eine Zuordnung von Aussagen
zu einzelnen Personen nicht möglich ist (vgl. Misoch 2015, S. 15--23).
Außerdem sollten sich die Befragten, sofern sie sich in den
Forschungsergebnissen \enquote{selbst wiedererkennen, durch die Art der
Darstellung nicht persönlich diskreditiert sehen}
(Przyborski/Wohlrab-Sahr 2021, S. 46). Mitunter wird deshalb überlegt,
die Interviewpartner*innen am Auswertungsprozess zu beteiligen. Dies
kann allerdings zu Lasten der Validität gehen: \enquote{Nicht immer
kommen wir in unseren Analysen zu Ergebnissen, die den Untersuchten
unmittelbar gefallen {[}\ldots{} Die{]} Differenz der Perspektiven ist
aber u. E. nicht dadurch auflösbar, dass man die beforschten Subjekte an
der Auswertung beteiligt oder ihnen die Auswertung und Verschriftlichung
gleichsam zur Freigabe vorlegt. Eine wissenschaftliche Bearbeitung
fördert andere Perspektiven zutage als Formen der Alltagsreflexion und
der Selbstbetrachtung. {[}\ldots{]} Auch die Anerkennung dieser
Differenz und der verantwortliche Umgang damit sind Ausdruck einer
forschungsethischen Haltung} (Przyborski/Wohlrab-Sahr 2021, S. 46).

\section{Konflikt: Probandenschutz versus Respekt geistigen
Eigentums}\label{konflikt-probandenschutz-versus-respekt-geistigen-eigentums}

Ein weiterer Wertekonflikt ergibt sich mit Blick auf die
Veröffentlichung von Daten: Gemäß den Regeln guter wissenschaftlicher
Praxis ist das geistige Eigentum anderer als solches kenntlich zu
machen. Demgegenüber steht der Schutz der Proband*innen, wonach
personenbezogene Informationen anonymisiert werden sollen (vgl. Misoch
2015, S. 15--23). Möchte man also aus Interviews zitieren, so darf man
dabei nicht die Namen der Interviewpartner*innen nennen. Eine Ausnahme
von dieser Regelung ist möglich, wenn ausdrücklich schriftlich
vereinbart wurde, dass die Proband*innen namentlich benannt werden
dürfen.\footnote{Das Online-Portal Scribbr bietet neben Anleitungen auch
  Musterdokumente für Interviewvorhaben. Vgl. \url{https://is.gd/jtFUqL}}
Ein Kompromiss zwischen den konfligierenden Werten wird zum Beispiel
erreicht, indem man die Interviewpartner*innen pseudonymisiert -- zum
Beispiel mit Bezeichnungen wie \enquote{Schüler*in 1, 2, 3 et cetera} /
\enquote{Kolleg*in 1, 2, 3 et cetera} / \enquote{Kund*in 1, 2, 3 et
cetera} -- und diese Pseudonyme nutzt, um die Urheberschaft von Zitaten
deutlich zu machen.

\section{Fazit}\label{fazit}

Das qualitative Interview ist eine gewinnbringende, aber keine
voraussetzungslose Methode. Forschungsethische Fragen werden auch in
qualitativen Forschungsdesigns \enquote{dauernd virulent, und zwar
deshalb, weil wir uns in eine Forschungsbeziehung mit einem menschlichen
Gegenüber begeben, und weil wir das, was dieses Gegenüber uns an
Persönlichem preisgibt, in Daten verwandeln, die einer
wissenschaftlichen Öffentlichkeit und bisweilen über diese hinaus
zugänglich gemacht werden sollen} (Przyborski/Wohlrab-Sahr 2021, S. 46,
im Original teilweise hervorgehoben). Wer qualitative Interviews führen
möchte, sollte sich daher mit forschungsethischen Fragen und den in
diesem Zusammenhang relevanten gesetzlichen Bestimmungen schon im
Vorfeld der Untersuchung auseinandersetzen, antizipieren, wo
gegebenenfalls Dilemmata auftreten und überlegen, wie die Proband*innen
vor Schaden geschützt werden können. Eine generelle Orientierung bieten
hierbei eine Reihe von Ethik-Kodizes, die von wissenschaftlichen
Organisationen erarbeitet und herausgegeben werden (vgl. Friedrichs
2022, S. 349). Diese enthalten oftmals zumindest auf allgemeiner Ebene
Stellungnahmen zu forschungsethischen Fragen, die sich auf den
Bibliothekskontext übertragen lassen.

\section{Literatur}\label{literatur}

Friedrichs, Jürgen: Forschungsethik. In: Baur, Nina; Blasius, Jörg
(Hrsg.): Handbuch Methoden der empirischen Sozialforschung. 3.
überarbeitete und erweiterte Aufl. Wiesbaden: Springer VS, 2022, S.
349--358. DOI: \url{https://doi.org/10.1007/978-3-658-37985-8}

Misoch, Sabina: Qualitative Interviews. Berlin: De Gruyter Oldenbourg,
2015. DOI: \url{https://doi.org/10.1515/9783110354614}

Przyborski, Aglaja; Wohlrab-Sahr, Monika: Qualitative Sozialforschung:
Ein Arbeitsbuch. 5. \linebreak überarbeitete und erweiterte Aufl. Berlin: De
Gruyter Oldenbourg, 2021. (Lehr- und Handbücher der Soziologie). DOI:
\url{https://doi.org/10.1515/9783110710663} Alexandra Corinna:
Bildungsverständnis und -ideal Öffentlicher Bibliotheken in Deutschland.
Berlin: Univ., Diss., 2024. DOI: \url{https://doi.org/10.18452/28067}

Seidman, Irving: Interviewing as Qualitative Research: A Guide for
Researchers in Education and the Social Sciences. Third Ed. New York:
Teachers College, Columbia University, 2006. URL:
\url{https://is.gd/qTlUPl}

Werner, Petra: Qualitative Befragungen. In: Umlauf, Konrad;
Fühles-Ubach, Simone; Seadle, Michael (Hrsg.): Handbuch Methoden der
Bibliotheks- und Informationswissenschaft: Biblio\-theks-,
Benutzerforschung, Informationsanalyse. Berlin: De Gruyter Saur, 2013,
S. 128--151. DOI: \url{https://doi.org/10.1515/9783110255546}

%autor
\begin{center}\rule{0.5\linewidth}{0.5pt}\end{center}

\textbf{Dr.~phil. Alexandra Roth} ist Lektorin bei der Münchner
Stadtbibliothek. Sie hat in Bibliothekswissenschaft promoviert und
arbeitet vor allem zu medienwissenschaftlichen und
bibliotheks\-pädagogischen Themen.

\end{document}